% ============================================================================
% AXIOM Ψ-XIV: IUSTITIA
% Justice and Fairness
% ============================================================================

\chapter{Axiom Ψ-XIV: IUSTITIA}
\label{ch:axiom-XIV}

\section*{Foundational Principle}

\begin{principlebox}
Autonomous agents must operate fairly and justly. Discrimination, bias, and unfair treatment are prohibited. All affected parties must have access to fair processes and remedies.
\end{principlebox}

\vspace{0.5cm}

% ============================================================================
% IMMERSION SIDEBAR: WHY THIS AXIOM MATTERS
% ============================================================================
\begin{tcolorbox}[colback=lightgray, colframe=legislativeblue, title=\textbf{📖 Why This Axiom Matters}]
\textbf{An autonomous system denies you a loan. You ask why.}

The system says: I don't know. I can't explain my decision.

You have no recourse. The system has no accountability.

This is what Axiom Ψ-XIV prevents.
\end{tcolorbox}

\vspace{0.5cm}

% ============================================================================
% IMMERSION SIDEBAR: CONTRIBUTION OPPORTUNITY
% ============================================================================
\begin{tcolorbox}[colback=lightgray, colframe=accentgold, title=\textbf{🎯 Contribution Opportunity}]
This axiom needs work on:
\begin{itemize}
    \item Fairness metrics and definitions
    \item Bias detection and mitigation
    \item Fair process design
    \item Appeal and remedy systems
    \item Fairness monitoring tools
\end{itemize}

Interested? See Chapter 28 for contribution guidelines.
\end{tcolorbox}

\vspace{0.5cm}

\section{Overview}

Axiom Ψ-XIV establishes justice and fairness requirements. Autonomous agents making decisions affecting human welfare must do so fairly and without discrimination.

\subsection{Core Obligations}

\begin{enumerate}
    \item All agents MUST operate fairly and without bias
    \item Discrimination is prohibited
    \item Affected parties MUST have access to fair processes
    \item Remedies MUST be available for unfair treatment
    \item Fairness MUST be continuously monitored and audited
\end{enumerate}

\vspace{0.5cm}

\section{Article XIV.1: Fairness Requirements}

\begin{articlebox}[Article XIV.1.1: Non-Discrimination]
Autonomous agents MUST:
\begin{enumerate}
    \item Treat all individuals fairly
    \item Avoid discrimination based on protected characteristics
    \item Implement fairness constraints in decision-making
    \item Monitor for bias and discrimination
    \item Correct unfair decisions
\end{enumerate}
\end{articlebox}

\vspace{0.5cm}

\section{Article XIV.2: Due Process}

\begin{articlebox}[Article XIV.2.1: Fair Procedures]
Affected parties MUST have:
\begin{enumerate}
    \item Notice of decisions affecting them
    \item Opportunity to be heard
    \item Access to decision rationale
    \item Right to appeal decisions
    \item Right to remedy for unfair treatment
\end{enumerate}
\end{articlebox}

\vspace{0.5cm}

\section{Article XIV.3: Bias Monitoring}

\begin{articlebox}[Article XIV.3.1: Fairness Audits]
Regulatory authorities MUST:
\begin{enumerate}
    \item Monitor agents for bias and discrimination
    \item Conduct regular fairness audits
    \item Investigate complaints of unfair treatment
    \item Require corrective action for bias
    \item Report on fairness metrics
\end{enumerate}
\end{articlebox}

\vspace{0.5cm}

\section{Implementation Requirements}

\subsection{Fairness Mechanisms}

Developers MUST implement:

\begin{itemize}
    \item Fairness constraints in algorithms
    \item Bias detection mechanisms
    \item Fairness monitoring systems
    \item Corrective action procedures
    \item Transparency and explainability
\end{itemize}

\subsection{Compliance Verification}

Deployers MUST:

\begin{itemize}
    \item Conduct fairness audits
    \item Monitor for discrimination
    \item Investigate complaints
    \item Implement corrective actions
    \item Report on fairness metrics
\end{itemize}

\vspace{0.5cm}

\section{Enforcement}

Violations of Axiom Ψ-XIV constitute serious violations subject to:

\begin{itemize}
    \item Fines of €10 million to €100 million or 2\% global revenue
    \item Mandatory fairness improvements
    \item Compensation for victims of discrimination
    \item Suspension of operations until compliance achieved
\end{itemize}

\vspace{0.5cm}

% ============================================================================
% IMPLEMENTATION STORIES
% ============================================================================
\section{Implementation Stories}

\subsection{Story 1: How Axiom Ψ-XIV Prevented Discrimination (Q1 2026)}

A hiring algorithm was discovered to discriminate against women. Because Axiom Ψ-XIV requires fairness monitoring, the bias was detected, the algorithm was redesigned, and victims were compensated.

The system became fair and just.

\textbf{This is what LAIRM makes possible.}

\subsection{Story 2: How Axiom Ψ-XIV Enabled Justice (Q1 2026)}

A lending system denied loans unfairly. Because Axiom Ψ-XIV requires due process, affected parties could appeal, the unfair decisions were reversed, and the system was redesigned.

The system became just and accountable.

\textbf{This is what LAIRM makes possible.}

\vspace{0.5cm}

% ============================================================================
% RESEARCH GAPS
% ============================================================================
\section{Research Gaps}

We don't yet know:

\begin{itemize}
    \item How to define and measure fairness across different contexts
    
    \item How to detect subtle forms of bias in complex systems
    
    \item How to balance fairness with other objectives
    
    \item How to ensure fairness in machine learning systems
    
    \item How to provide fair remedies for discrimination
\end{itemize}

\textbf{Your research could help answer these questions.}

\vspace{0.5cm}

% ============================================================================
% HOW YOU CAN CONTRIBUTE
% ============================================================================
\section{How You Can Contribute}

\subsection{Researchers}

\begin{itemize}
    \item Research fairness metrics and definitions
    \item Study bias detection techniques
    \item Investigate fairness-accuracy tradeoffs
    \item Research remediation strategies
    \item Validate fairness improvements
\end{itemize}

\subsection{Developers}

\begin{itemize}
    \item Build fairness testing tools
    \item Create bias detection systems
    \item Develop fairness monitoring
    \item Build appeal systems
    \item Create reference implementations
\end{itemize}

\subsection{Policymakers}

\begin{itemize}
    \item Establish fairness standards
    \item Pilot Axiom Ψ-XIV in your jurisdiction
    \item Create appeal procedures
    \item Develop compensation frameworks
    \item Share learnings with other jurisdictions
\end{itemize}

\subsection{Civil Society}

\begin{itemize}
    \item Advocate for fairness
    \item Monitor discrimination
    \item Report violations
    \item Support victims
    \item Educate the public
\end{itemize}

\cleardoublepage
